\addchap{Agradecimientos}

La creación de esta nueva plantilla para \LaTeX{} es el resultado del esfuerzo acumulado de numerosos compañeros de la Escuela. En primer lugar, es obligado reconocer la labor del Prof. Payán, quien, de forma desinteresada y partiendo de su experiencia en la edición de libros docentes, dio vida a la versión original de este formato; un diseño que se ha convertido en un símbolo de la imagen institucional de la ETSi. Del mismo modo, agradezco la dedicación de los profesores Juan José Murillo y Germán Madinabeitia, piezas clave en el mantenimiento de la plantilla original durante años.

Por otra parte, las portadas incluidas tienen su origen en los diseños realizados por el Prof. Fernando García García, de la Facultad de Bellas Artes, para la sección de publicaciones de nuestra Escuela. Nuestra Escuela le agradece que pusiera su arte y su trabajo, de forma gratuita, a nuestra disposición.

Finalmente, expreso mi gratitud al Profesor Antonio Estepa, tanto por su confianza para encomendarme esta labor como por las valiosas recomendaciones de redacción que ha aportado al resumen y al primer capítulo de este documento.

{\flushleft{\hfill \textit{Vicente Jesús Mayor Gallego} \\
\hfill \textit{Profesor del Departamento de Ingeniería Telemática}}}%
{\flushleft{\hfill \emph{Sevilla, 2026}}}%

\addchap{Resumen} 

El resumen de un Trabajo Fin de Grado o Máster debe ofrecer al lector una visión clara y concisa de \textbf{todo} el contenido del proyecto. No es ni una mera introducción ni una tabla de contenidos. Por ello, debería incluir, al menos, el contexto y problema abordados, el objetivo del trabajo, el método seguido, los principales resultados obtenidos, y principales conclusiones o implicaciones de los resultados. Debe ser autocontenido (p.ej., en lugar de escribir ``en la sección 3 se muestran los resultados'', se resumen los principales resultados de la sección). \textbf{El resumen tendrá una extensión entre 150 y 500 palabras}, suficiente para que cualquier lector pueda captar el propósito, desarrollo e importancia del trabajo, y decidir si es de su interés la lectura del proyecto o alguna parte del mismo.  El lenguaje empleado ha de ser formal, directo y fácil de seguir, evitando tecnicismos o detalles innecesarios que dificulten la comprensión. 

Ejemplo de un resumen corto de un proyecto del MIT: \textit{Este proyecto final presenta un sistema robótico concebido bajo el concepto de ``Legos industriales'' modulares, con el objetivo de aprovechar de una forma innovadora las ventajas de la modularidad. El sistema está compuesto por módulos capaces de autoensamblarse dentro de estructuras reticulares previamente definidas. Las principales aportaciones son: (1) el diseño de un módulo compacto que integra actuación, percepción del entorno y conexión magnética; (2) algoritmos de control que permiten a los módulos unirse de manera autónoma y formar la configuración deseada; y (3) un procedimiento de detección y corrección de errores que aumenta la robustez y la precisión del ensamblaje. Los experimentos realizados con un prototipo muestran que el sistema puede ensamblar retículas de hasta 12 módulos con una tasa de éxito del 90\% en condiciones normales, y que es capaz de recuperarse ante determinados fallos. Estos resultados indican que el sistema propuesto tiene un gran potencial para el desarrollo de sistemas escalables y reconfigurables, con aplicaciones que van desde la manufactura adaptable hasta las estructuras desplegables.}

Cualquier persona que lea tu resumen debería ser capaz de responder a las siguientes preguntas: (a) ¿cuál es el tema o problema u objetivo tratado en el proyecto y por qué es importante?; (b) ¿qué métodos o técnicas o soluciones computacionales has utilizado para obtener los resultados de tu proyecto?; (c) ¿cuáles son los principales resultados que has obtenido? (cuantitativamente, si fuera posible); (d) ¿qué implican los resultados o cómo avanzan el campo de aplicación tratado?


\addchap{Abstract} 

Una traducción al inglés del resumen anterior. Incluye el título del proyecto también en el resumen en inglés. 
